\chapter{Four harder tests}

Everything so far used tidy, ready-made records: one global temperature per year, one for the United States. A fair
reader could object that I never went back to the thermometers the study used, that I ignored El Niño, that I treated
the global record as if it were known exactly, and that a two-reservoir model is not a climate model. Each objection
can be tested, so I tested them. The hands-on guide
\href{https://ahmadammar.github.io/teaching/climate-data-python/}{Climate data in Python} walks through every step in
code.

\section{Back to the thermometers}

NOAA keeps the monthly records of about 27,000 weather stations in two versions: as reported, and after
\emph{homogenisation}, the correction for station moves, new instruments and changed reading times described earlier under
``Old thermometers need correcting''. The station identifiers are the same as in the daily archive the study used, so I could rebuild the study's
own selection (stations that measured daily maximum and minimum from 1900 or earlier to 2020 or later) and find
983 of them with a usable monthly record. For those stations, for all land stations and for the United States, I
turned each station's temperatures into departures from its own 1951--1980 average and combined them in two ways:
every station counting once, which is what the study's method does, and by area, averaging stations inside
5-degree boxes first and weighting each box by its size.

\begin{figure}[H]
\centering
\includegraphics[width=\linewidth]{stations_trends.png}
\caption{Warming rate in each of the paper's four periods (bars, oldest light grey to newest red), for twelve ways of
building the record: raw or corrected (homogenised) data; every station counted equally or weighted by area; and the
study's own stations, all land stations, or the contiguous United States. Data: NOAA GHCN-Monthly version 4.}\label{fig:stations_trends}
\end{figure}

\textbf{How to read Figure~\ref{fig:stations_trends}.} Each group of four bars is one way of building the record. Within
a group, the yellow bar is 1899--1940 and the red bar is 1983--2024; the question is which is taller.

The result explains the study's finding exactly. Take the study's own stations, uncorrected, every station counted
once: the early and recent periods come out almost the same, 0.20 and 0.21~°C per decade. That is close to what the
study reported. Now change one thing at a time. Correct the data, and the recent period rises to 0.30. Weight by area
instead, and the early period falls to 0.12, because the United States no longer stands in for the world. Do both,
as every serious temperature record does, and recent warming is 0.36~°C per decade, nearly three times the early rate
of 0.13. The study's ordering was made by its data choices, not by the climate.

\begin{figure}[H]
\centering
\includegraphics[width=\linewidth]{stations_maps.png}
\caption{Corrected stations with at least 30 complete years in each period. Red: United States; blue: elsewhere.}\label{fig:stations_maps}
\end{figure}

\textbf{How to read Figure~\ref{fig:stations_maps}.} Each dot is a station that has a long record in that period.
Before 1940 the long records are almost all in the United States and Europe; most of Africa, South America and Asia is
empty. An average that counts stations rather than area is, for those years, an average of the United States.

\section{El Niño}

Every few years the tropical Pacific releases a pulse of heat into the air, an El Niño, and the whole planet is
warmer for a year; its opposite, La Niña, cools it. If one period happened to begin in a La Niña and end in an El
Niño, its trend would look steeper than the climate's. I took the standard index of these swings, the sea
temperature of the central Pacific since 1854, and removed its share from the global record.

\begin{figure}[H]
\centering
\includegraphics[width=\linewidth]{enso.png}
\caption{Top: global temperature as measured (grey), with El Niño and La Niña removed (blue), and the response of the
two-reservoir model to all forcings (red). Bottom: the Niño~3.4 index, three months earlier; red bars El Niño, blue
La Niña. Data: HadCRUT5; NOAA ERSSTv5 via KNMI Climate Explorer.}\label{fig:enso}
\end{figure}

\textbf{How to read Figure~\ref{fig:enso}.} The blue line is smoother than the grey one: the spikes of 1878, 1998 and
2016 are gone. Yet the long shape is unchanged. El Niño and La Niña explain part of the year-to-year jitter (the share
of the variation explained rises from 87\% to 92\%), but they shift no period's trend by more than 0.02~°C per decade.
They are noise around the trend, not its cause.

\section{How sure are the records?}

Nobody knows the temperature of 1880 exactly. The Met Office therefore publishes 200 equally likely versions of the
global record, which agree where measurements were plentiful and spread where they were thin. I computed every
period's trend on all 200 versions, adding the uncertainty from the unmeasured parts of the globe and the uncertainty
of fitting a line.

\begin{figure}[H]
\centering
\includegraphics[width=0.85\linewidth]{uncertainty.png}
\caption{The warming rate of each period across 200 versions of the global record (HadCRUT5), including coverage and
fitting uncertainty. The width of each shape shows how many results fall at that value; the black line spans 95\% of
them.}\label{fig:uncertainty}
\end{figure}

\textbf{How to read Figure~\ref{fig:uncertainty}.} The older periods have wider shapes, because the early record is less
certain. The early period (1899--1940) lies between 0.05 and 0.16~°C per decade; the recent one (1983--2024) between 0.19
and 0.24. The two ranges do not touch: in all of the 200 versions, the recent period warmed faster.

\section{What full climate models say}

My two-reservoir model is a sketch. Full climate models divide the atmosphere and oceans into boxes about 100~km
wide and solve the physics of air, water and heat in each, step by step. Modelling centres around the world run them
with the same historical forcing for a shared experiment, CMIP6. Each model also makes its own weather, including its
own El Niños at random times, so no single model should match the real wiggles, but together they show what the
forcing, plus chance, produces. I read the global temperature of 34 of them straight from their public cloud archive.

\begin{figure}[H]
\centering
\includegraphics[width=\linewidth]{cmip6.png}
\caption{Left: global temperature from 34 climate models (red line, middle model; pink band, the middle 90\%) and as
measured (black, HadCRUT5), relative to 1850--1900. Right: the warming rate of each period in the models (light bar,
full range; dark bar, middle 90\%) and as measured (black dot). Data: CMIP6 historical and SSP2-4.5 runs via the
Pangeo cloud archive.}\label{fig:cmip6}
\end{figure}

\textbf{How to read Figure~\ref{fig:cmip6}.} On the left, the black line stays mostly inside the pink band: models
given only the forcing reproduce the slow rise to 1940, the flat middle of the century and the steep climb since. On
the right, each bar is the spread of the models' rates for one period, and the dot is what was measured. In all 34
models, 1983--2024 warms faster than any earlier period. The one place the measurements stand out is 1910--1945: the
real world warmed faster than 33 of the 34 models. That is the unforced, natural part of the early warming that the
previous chapter's ``unexplained'' bar showed (Figure~\ref{fig:attribution}), now seen in full models too.


\section{What the tests add}

None of the four tests changes the answer of the previous chapters, and one of them explains where the study's answer
came from. Its central result can be reproduced, but only from uncorrected thermometers counted one by one, most of
them in one country. Correct the thermometers or count by area, and the recent decades warm fastest; account for El
Niño, for the uncertainty of the record, or compare with full climate models, and they still do.
